\section{Canonical upstream reduction}\label{sec:upstream}
\subsection{A certified parameter interval}
\begin{proposition}[Finite certificates for the initial interval]\label{prop:R1}
With the exact rationals \(L,U\) of Theorem~\ref{thm:main},
\begin{equation}\label{eq:R1}
 L<\St\le U.
\end{equation}
\end{proposition}
\begin{proof}
The lower certificate reconstructs six rational orthogonal projectors on
\(\CC^{499}\) and a nonzero rational Schmidt vector. The exact circle
identities for all 498 adjacent block parameters make the blocks projective.
Dividing their direct Bell contraction by the positive squared norm gives a
rational value \(v\) satisfying
\[
 0.2508753845139765354\le v
 <0.2508753845139765355.
\]
It is an actual finite tensor value, so \(L<v\le\St\).

For the upper bound, let \(\mathbf w\) be the column of 244 reduced words of
total length at most four in the six involutions of
\eqref{eq:involutions}. Words are reduced using adjacent-square cancellation
within each party and commutation across parties only. The upper certificate
constructs a rational positive semidefinite \(244\times244\) matrix \(Y\)
and checks the noncommutative polynomial identity
\begin{equation}\label{eq:upperidentity}
 \frac{2501750771}{500000000}I-\mathsf B
       =\mathbf w^*Y\mathbf w.
\end{equation}
The positivity check consists of 183 positive exact leading principal
determinants of five integer blocks, followed by a positive congruence.
Expansion of all 59,536 ordered word products has zero residual at all
8,452 coefficient positions. Evaluating the identity in any finite tensor
strategy and any vector gives a nonnegative right-hand expectation.
Equation~\eqref{eq:involutions} therefore gives
\[
 \Bell\le \frac14\left(\frac{2501750771}{500000000}-4\right)I
       =\frac{501750771}{2000000000}I=UI.
\]
This constant is uniform in the finite dimensions, so it bounds their
supremum. Appendix~\ref{app:bounds} specifies the rational inputs, block
matrices, reconstruction and acceptance predicates. No numerical solver
output, hierarchy convergence theorem or symmetry-completeness assertion is
needed for this operator implication.
\end{proof}

\subsection{Precisely restricted external premises}
Put
\begin{equation}\label{eq:coefficients}
 w(c)=\sqrt{1-c^2},\qquad b(c)=\frac{w(c)}2,\qquad
 d(x,y)=xy+\frac{x-y}{2}-1.
\end{equation}
For a finite path \(\gamma=(\gamma_0,\ldots,\gamma_m)\) with
\(\gamma_0,\gamma_m\in\{-1,1\}\) and the remaining labels in
\([-1,1]\), define
\begin{equation}\label{eq:J}
 (J_m(\gamma))_{ii}=d(\gamma_{i-1},\gamma_i),\qquad
 (J_m(\gamma))_{i,i+1}=b(\gamma_i),
 \qquad q^*=\sup_{m,\gamma}\lambda_{\max}(J_m(\gamma)).
\end{equation}
The matrices are real symmetric; unspecified entries vanish.

\begin{premise}[Coladangelo v1]\label{prem:Coladangelo}
For the functional \eqref{eq:bell} and the number \(q^*\) in
\eqref{eq:J}, we use exactly the following three source interfaces:
\begin{enumerate}
\item Theorem 7.1 of~\cite{Coladangelo2026v1} (p.~24, (33)--(34):
there is a finite, nonnegative, concave function
\(\mathfrak h:[-1,1]\to[0,\infty)\) such that, on the closed intervals,
\begin{equation}\label{eq:storage}
 \mathfrak h(y)+d(x,y)+\mathfrak h(-x)\le q^*,\qquad
 \mathfrak h(z)\mathfrak h(-z)\ge b(z)^2.
\end{equation}
\item The finite-dimensional projective tensor restriction of Theorem 7.2
of that source (p.~24): every normalized strategy in
Definition~\ref{def:St} has value at most \(q^*\).
\item Proposition 10.9 of that source (pp.~70--71): there are finite paths
\(1=\gamma^{(n)}_0\ge\cdots\ge\gamma^{(n)}_{m_n}=-1\)
whose largest Jacobi eigenvalues tend to \(q^*\).
\end{enumerate}
\end{premise}
These are three consumed interfaces from one source, with shared proof
ancestry. In particular the upper-bound theorem uses the scalar-certificate
theory. They are not three independent external confirmations. No global
differentiability or strict concavity of \(\mathfrak h\), and no classification of its
entire equality graph, is part of the premise. The source's broader
representation statements are not imported.

\begin{proposition}[Finite realization and value identification]\label{prop:qstar}
Every largest eigenvalue in \eqref{eq:J} is realized by a finite-dimensional
projective tensor strategy for \eqref{eq:bell}. Consequently \(q^*=\St\).
Moreover for every bilateral sequence \(c_j\in[-1,1]\), the bounded
self-adjoint Jacobi operator
\begin{equation}\label{eq:H}
 (H(c)f)_j=b(c_j)f_{j-1}+d(c_j,c_{j+1})f_j+b(c_{j+1})f_{j+1}
\end{equation}
satisfies \(H(c)\le\St I\).
\end{proposition}
\begin{proof}
Appendix~\ref{app:realization} constructs the six projectors explicitly by
alternating two-dimensional blocks. On
\(\sum_i v_i e_i\otimes e_i\), the Bell expectation is exactly
\(v^TJ_m(\gamma)v\). A real unit top eigenvector therefore proves
\(q^*\le\St\). External premise~\ref{prem:Coladangelo}(2) proves the reverse
inequality in precisely the same finite projective tensor class.

On \([-1,1]^2\), \(-3\le d\le0\), since its four corner values are
\(0,-1,-3,0\) and bilinear interpolation has nonnegative weights. Also
\(0\le b\le1/2\), hence \(\|H(c)\|\le4\). For any finitely supported
real \(f\), retain all labels incident on its support, add a zero-amplitude
vertex at either end, and choose the two new outer labels in \(\{-1,1\}\).
The resulting finite path has the same Rayleigh numerator. Its value is at
most \(q^*\|f\|^2\). Density and boundedness extend the inequality to
\(\ell^2(\ZZ)\); real and imaginary parts give the complex version. This
comparison permits arbitrary finite label changes, with no monotonicity
constraint.
\end{proof}
